% =====================================================================
%  CHƯƠNG 4 — OCR & MÀU BIỂN   | Tuần 3-4 | Phụ trách: Nhật (OCR) + Đức (màu/ánh sáng)
% =====================================================================
\chapter{Nhận Dạng Ký Tự và Màu Nền Biển Số}
\label{ch:ocr-maubien}

\section{Tiền Xử Lý Vùng Biển Số}
\wip{Cắt vùng biển từ bộ phát hiện; hiệu chỉnh nghiêng (deskew); tách 2 hàng cho biển xe máy.}

\section{Nhận Dạng Ký Tự (OCR)}
\subsection{Lựa Chọn Công Cụ OCR}
\wip{So sánh PaddleOCR vs EasyOCR \autocite{paddleocr2025ppocrv5,moussaoui2024yolov8ocr}; lý do chọn.}
\subsection{Tinh Chỉnh và Hậu Xử Lý}
\wip{Fine-tune trên biển VN; luật hậu xử lý (định dạng biển VN, sửa nhầm O/0, B/8...).}
\subsection{Đánh Giá OCR}
\wip{Độ chính xác mức biển (plate-level) + mức ký tự / CER. Bảng kết quả.}

\section{Nhận Diện Màu Nền Biển Số}
\subsection{Xử Lý Ánh Sáng}
\wip{CLAHE + chuẩn hóa ánh sáng; xử lý điều kiện thiếu/thừa sáng.}
\subsection{Phân Loại Màu bằng Không Gian HSV}
\wip{Ngưỡng HSV cho 4 màu (trắng/vàng/xanh/đỏ) $\rightarrow$ nhóm xe cá nhân/kinh doanh/nhà nước.}
\subsection{Đánh Giá Màu Biển}
\wip{Accuracy theo màu + ma trận nhầm lẫn; đánh giá đa điều kiện ánh sáng.}

\section{Thảo Luận}
\wip{Nguồn lỗi chính (mờ, chói, nghiêng); giới hạn và hướng cải thiện.}
